\section{Wings are tails: Lee's bound inside the density factor}
\label{sec:svilee}

Why should a wing slope be capped at all?  Because it is a moment
barometer.  Lee's formula \eqref{eq:lee} ties the asymptotic slope of total
variance to the critical moment order of the gross return: on the right
wing $\beta_R=\Psi(r_+^*)$, where $1+r_+^*$ is the largest finite moment of
$Y$ and $\Psi(r)=2-4(\sqrt{r^2+r}-r)$.  As the law loses its moments,
$r_+^*\downarrow0$ and $\Psi\uparrow2$: a wing steeper than $2$ is the tail
of no probability law at all.  For LQD the correspondence was a closed
formula in the endpoint speeds (\cref{prop:leeclosed}); for a family whose
wings are exactly linear it is a sharp, checkable inequality on two
parameters, $\max(\beta_L,\beta_R)\le2$.

The cleanest way to watch the bound bite is to follow it into the function
that tests for arbitrage.  The Durrleman factor \eqref{eq:durrleman} mixes
$w$, $w'$, and $w''$; in the wings of the family \eqref{eq:svi} the mixture
simplifies to a limit that depends on the slope alone.

\begin{proposition}[The wing slope sets the sign of the tail density]
\label{prop:gtail}
For a slice \eqref{eq:svi} with $w>0$,
\begin{equation}
\lim_{k\to-\infty}g_{\rm D}(k)=\frac{4-\beta_L^2}{16},\qquad
\lim_{k\to+\infty}g_{\rm D}(k)=\frac{4-\beta_R^2}{16}.
\label{eq:gtail}
\end{equation}
Hence for $\beta<2$ the corresponding tail of $g_{\rm D}$ is eventually
positive and for $\beta>2$ eventually negative.  At $\beta_R=2$ the limit
vanishes and the sign falls to the next order: the right wing is
$w(k)=2k+(a_S-2m_S)+O(k^{-1})$, and
\begin{equation}
g_{\rm D}(k)=\frac{(a_S-2m_S)-2}{4k}+O(k^{-2}),
\label{eq:gboundary}
\end{equation}
so the boundary slice carries negative tail density whenever
$a_S-2m_S<2$.  The mirror statement holds on the left with the intercept
$a_S+2m_S$.
\end{proposition}

\begin{proof}
In either wing $w\sim\beta|k|$, $|w'|\to\beta$, and $w''\to0$ by
\eqref{eq:svideriv}.  The first bracket of \eqref{eq:durrleman} tends to
$1-\tfrac12=\tfrac12$, so its square tends to $\tfrac14$; the middle term
tends to $\beta^2/16$ because $1/w\to0$; the last term vanishes.  This
gives \eqref{eq:gtail}.  For the boundary order, expand \eqref{eq:svi} on
the right with $b_S(1+\rho_S)=2$:
\[
w(k)=2k+(a_S-2m_S)+\frac{b_Ss_S^2}{2k}+O(k^{-2}),
\qquad
w'(k)=2+O(k^{-2}).
\]
Then $kw'/(2w)=\tfrac12\{1-(a_S-2m_S)/(2k)\}+O(k^{-2})$, so the first
bracket squares to $\tfrac14+(a_S-2m_S)/(4k)+O(k^{-2})$; the middle term is
$\tfrac14\,w'^2\{1/w+\tfrac14\}=\tfrac14+1/(2k)+O(k^{-2})$ because
$w'^2\to4$ and $1/w=1/(2k)+O(k^{-2})$; and $w''=O(k^{-3})$.  Subtracting
gives \eqref{eq:gboundary}.
\end{proof}

The proposition says two distinct things, and the second is easy to miss.
Below the bound, the tail of the implied density is eventually positive ---
the wing condition genuinely protects the far tail.  \emph{At} the bound,
it protects nothing: the limit is zero, and the $1/k$ term decides.  A
concrete slice makes the boundary failure explicit.  Take
$(a_S,b_S,\rho_S,m_S,s_S)=(0.04,\,2,\,0,\,0,\,0.2)$: its minimum total
variance is $0.44$, comfortably positive, and both wings sit exactly at
$\beta=2$, so any constraint of the form $\beta\le2$ is satisfied with
equality and charges nothing.  Yet $a_S-2m_S=0.04<2$, and by
\eqref{eq:gboundary} the density is negative throughout the far tail ---
$g_{\rm D}(10)=\MacSviBoundaryGTen$ on this slice.  A constraint written
with Lee's constant as its right-hand side therefore admits laws that are
not laws.

The reference implementation draws the consequence as a convention, stated
here once: the enforced cap is strictly inside the bound,
\begin{equation}
\max(\beta_L,\beta_R)\;\le\;\beta_{\max},\qquad
\beta_{\max}=\MacSviBetaMax<2 .
\label{eq:leecap}
\end{equation}
What does the buffer exclude?  Inverting $\Psi$ gives the \emph{moment
budget} at slope $\beta$,
\begin{equation}
r^*=\Psi^{-1}(\beta)=\frac{(2-\beta)^2}{8\beta},
\label{eq:momentmap}
\end{equation}
the number of moments beyond the first that a wing of slope $\beta$ still
permits.  At the cap, $r^*=\MacSviPstarCap$: the laws excluded by the
buffer are those whose entire moment budget beyond the mean is smaller
than two parts in ten thousand.  The cap costs essentially nothing in
distributional scope and restores the implication ``wing-admissible
$\Rightarrow$ tail density eventually positive.''

The map \eqref{eq:momentmap} also says when a fitted wing may be read as a
moment estimate, because its conditioning is extreme at both ends.
Differentiating, $\dd r^*/\dd\beta=-(4-\beta^2)/(8\beta^2)$, which behaves
as $-1/(2\beta^2)$ for small $\beta$: a shallow wing pins essentially no
upper bound on the moments, and small changes in the fitted slope move the
inferred moment violently.  As $\beta\to2$ the budget collapses to zero.
A fitted $\beta$ is a candidate moment reading only in the middle of the
range --- and only for a slice that is actually a law, which is the
subject of the next section.  \Cref{fig:svilee} draws all three parts of
the story.

\begin{figure}[htbp]
\centering
\paperfig{\textwidth}{fig_svi_lee.pdf}
\caption{Lee's bound inside the density factor.  (a)~The tail limit
\eqref{eq:gtail} is a downward parabola crossing zero exactly at
$\beta=2$; the dashed line is the working cap $\beta_{\max}=\MacSviBetaMax$.
(b)~The boundary slice of the text: at $\beta_R=2$ (orange) the tail
density is negative at every plotted strike, approaching zero from below
as \eqref{eq:gboundary} prescribes; the same slice with its wings at the
cap (blue) turns positive and settles at the limit
$(4-\beta_{\max}^2)/16$ (dotted).  (c)~The moment budget
\eqref{eq:momentmap} on a log scale: it explodes as $\beta\to0$ and
collapses at the cap (marked point, $r^*=\MacSviPstarCap$).}
\label{fig:svilee}
\end{figure}
